%% This is based on the leaflet.tex file present in the texlive-leaflet module
%% 
\def\filename{fedora-workstation-flyer.tex}
\documentclass[
%%notumble,
%%nofoldmark,
%%dvipdfm,
%%portrait,
%%titlepage,
%%nocombine,
%%a3paper,
%%debug,
%%nospecialtricks,
%%draft,
]{leaflet}


\renewcommand*\foldmarkrule{.3mm}
\renewcommand*\foldmarklength{5mm}

\usepackage[T1]{fontenc}
\usepackage{textcomp}
\usepackage{mathptmx}
\usepackage[scaled=0.9]{helvet}
\makeatletter
\def\ptmTeX{T\kern-.1667em\lower.5ex\hbox{E}\kern-.075emX\@}
\DeclareRobustCommand{\ptmLaTeX}{L\kern-.3em
        {\setbox0\hbox{T}%
         %\vb@xt@ % :-)
         \vbox to\ht0{\hbox{%
                            \csname S@\f@size\endcsname
                            \fontsize\sf@size\z@
                            \math@fontsfalse\selectfont
                            A}%
                      \vss}%
        }%
        \kern-.12em
        \ptmTeX}
\makeatother
\let\TeX=\ptmTeX
\let\LaTeX=\ptmLaTeX
\usepackage{shortvrb}
\MakeShortVerb{\|}
\usepackage{url}
\usepackage{graphicx}
%%\usepackage[default]{cantarell} %% Use option ``defaultsans'' to use cantarell as sans serif only
\usepackage[default]{comfortaa} %% looks better
\usepackage[T1]{fontenc}
\usepackage[dvipsnames,usenames]{color}
\definecolor{LIGHTGRAY}{gray}{.9}

%%%%\renewcommand{\descfont}{\normalfont}
\newcommand\Lpack[1]{\textsf{#1}}
\newcommand\Lclass[1]{\textsf{#1}}
\newcommand\Lopt[1]{\texttt{#1}}
\newcommand\Lprog[1]{\textit{#1}}

\newcommand*\defaultmarker{\textsuperscript\textasteriskcentered}

%%\title{Fedora workstation}
\title{\includegraphics{workstation-logo.png}}
\date{December 10, 2014}
\author{}

\CutLine*{1}% Dotted line without scissors
\CutLine{6}%  Dotted line with scissors

%% Use better logo for the workstation
%%\AddToBackground{1}{%  Background of a small page
%%  \put(20,530){\includegraphics{workstation-logo.png}}}

\AddToBackground{5}{%  Background of a small page
  \put(0,0){\textcolor{Cerulean}{\rule{\paperwidth}{\paperheight}}}}

\AddToBackground*{2}{% Background of a large page
  \put(\LenToUnit{.5\paperwidth},\LenToUnit{.5\paperheight}){%
    \makebox(0,0)[c]{%
      \resizebox{.9\paperwidth}{!}{\rotatebox{35.26}{%
        \textsf{\textbf{\textcolor{LIGHTGRAY}{BACKGROUND}}}}}}}}

\begin{document}

\maketitle
\thispagestyle{empty}

%%\LARGE

%%\tableofcontents

\section{Overview}

The document class \Lclass{leaflet} creates a document of (up to) six
small pages in portrait orientation, arranged physically on two
``normal-size'' pages. The target page sizes supported by the standard
\LaTeX{}  \Lclass{article} are available, plus |a3paper|. Printing these
to both sides of a sheet and folding appropriately will yield a six-page
leaflet.

%% TEST: These commands are no longer disabled!
\iffalse
\reversemarginpar
\marginpar[XXX]{YYY}
\fi
%% TEST: These commands are disabled!
\onecolumn
\twocolumn[WWW]
%% end of TEST

\section{Requirements}

Using the \Lclass{leaflet} class requires that the final
document is created in PostScript or PDF format, using
\begin{itemize}
  \item \TeX{} and \Lprog{dvips}, or
  \item pdf\TeX{}, or
  \item V\TeX{} in PS or PDF mode.
\end{itemize}
(Some other drivers supported by standard \LaTeX{} work as well.)

The non-standard macro package \Lpack{everyshi} \cite{cit:everyshi} is
used by the \Lclass{leaflet} class.

\section{Features}

Basically the \Lclass{leaflet} class provides the same features as the
standard \Lclass{article} class. There are, however, a number of
differences and restrictions, as well as some additional facilities and
peculiarities:

\begin{itemize}

\item
The sectioning level |\part| is not available.
The other sectioning levels are not numbered by default.

\item
References to the page where floating objects are located may come out
wrong (this includes |\pageref| as well as |\listof...| commends).

\item
Section headers are typeset in a smaller font size than in the
standard classes.

\item
You may use list-like environments just as in the standard classes.
The left margins have been adjusted to work well with the
\Lopt{a4paper} and \Lopt{letterpaper} class options.
With other target page sizes, you'll have to adjust them.

Here's a small demo:
\begin{description}
\item[Uncle Meat] First entry in a description environment.
\item[King Kong] Second entry.
    \begin{itemize}
    \item First entry in an itemize environment.
        \begin{enumerate}
        \item First entry in an enumerate environment.
        \item Second entry.
            \begin{enumerate}
            \item First entry in an enumerate environment.
                \begin{enumerate}
                \item First entry in an enumerate environment.
                \item Second entry.
                \end{enumerate}
            \item Second entry.
                \begin{itemize}
                \item First entry in an itemize environment.
                \item Second entry.
                \end{itemize}
            \item Another entry.
            \end{enumerate}
        \item Another entry.
        \end{enumerate}
    \item Second entry.
    \item Another entry.
    \end{itemize}
\item[Frunobulax] Another entry.
\end{description}

\item
Marginal notes are pointless on the given page size and
are disabled.

\item
Two-column typesetting is not supported for the same reason.

\item
By default, there are no page headers, page footers or page numbers,
nor is there any space reserved for these.

However, you can restore them, if you like.
To do so, use |\pagestyle| as with the standard classes,
and |\setlength| to adjust the corresponding parameters (like
|\headheight|).
At last, you have to call the new macro \par
|  \setmargins{top}{bottom}{left}{right}|.

\item
Paragraphs are separated by vertical space; the first line
of a paragraph is not indented by default.

\item
By default, all paragraphs are typeset as if you had
specified |\sloppy| in the document preamble.

\item
A small folding mark is created between the second and the
third page.

\item
The macro |\CutLine| draws a vertical dotted line with
scissor symbols between the page indicated by its argument and the
preceding one.
The starred version omits the scissors symbols.

\item
In case the text does not fit on six pages, a warning (or error,
depending on some class option, see below) will be issued.

\item
To add some background picture to individual pages, you can use
|\AddToBackground| commands. Its first argument specifies the page,
the second one the picture commands.
The starred version puts the picture on the combined pages.

\end{itemize}

\section{Customization}

The typeface to be used for the section headings is given by the macro
|\sectfont|, and the typeface to be used for the labels of the
|description| environment is given by |\descfont|.
Both macros default to |\bfseries| and can be changed using
|\renewcommand*|.

The horizontal and vertical and margins of the (small) pages
default to 8\,mm and 11\,mm, respectively, and can be changed
using |\setmargins|, as explained above.
This may be useful, if the printing engine exhibits larger unprintable
margins.

The macros |\foldmarkrule| and |\foldmarklength| determine the stroke
width and the length of the fold mark, which is printed between the
second and the third page.
They default to 0.4\,pt and 2\,mm, respectively, and can be changed
using |\renewcommand*| (\emph{not} |\setlength|!).
See also the class options \Lopt{foldmark} and \Lopt{nofoldmark}.

\section{Class options}

Default options are marked with an asterisk:
\begin{description}
\item[\Lopt{tumble}{\defaultmarker}, \Lopt{notumble}]
  By default, the contents of the back side of the final sheet is
  printed upside down.
  The option \Lopt{notumble} suppresses that.
  Doing so may be necessary to suit the behavior of certain printing
  engines.
  Specifying \Lopt[notumble] may also be useful during the writing of
  a document, to enable proof-reading on the screen.
\item[\Lopt{frontside}, \Lopt{backside}, \Lopt{bothsides}{\defaultmarker}]
  These options control whether only the front page, the back page or
  both pages of the final sheet are to be be created.
  Thus, you can create separate files for the front an back side of
  the sheet.
\item[\Lopt{foldmark}{\defaultmarker}, \Lopt{nofoldmark}]
  These options specifiy whether or not a fold mark is to be printed.
\item[\Lopt{combine}{\defaultmarker}, \Lopt{nocombine}]
  These options specify whether the (small) pages should be output
  combined on a (large) target page (\Lopt{combine}) or as individual
  pages (\Lopt{nocombine}).

  At the same time, the determine behaviour in case the text does not
  fit on six (small) pages.
  By default (\Lopt{combine}), an error is raised---and the surplus pages
  will be gobbled.
  Otherwise (\Lopt{nocombine}), just a warning will be issued;
  shortening the text appropriately is left to the user.
\item[\Lopt{twopart}, \Lopt{notwopart}{\defaultmarker}]
  Allows the typesetting of a four page leaflet (first part) and a two page
  detachable form (second part), for fill-in forms, questionnaires, applications,
  etc.
\end{description}

Other options are passed to the \Lclass{article} class.

\section{Changes over version 0.3}

The present release of the \Lclass{leaflet} class differs basically
from its predecessor, version~0.3, which had been developed originally
by J\"urgen Schlegelmilch.

The main change is, that no more post-processing is required to
arrange the pages on the sheet.
Furthermore, the overall layout has been changed slightly to suit the
small page size better.
In general, documents that were written for version~0.3 will exhibit
different line and page breaks when typeset using the new version of
this document class.

\begin{thebibliography}{000}
\bibitem{cit:latex-man}
  \textsc{L.\,Lamport}: \LaTeX. A Document Preparation System.
  \textit{User's Guide And Reference Manual.} Second Edition. 1994.
\bibitem{cit:everyshi}
  \textsc{M.\,Schr\"oder}: The \Lpack{everyshi} package. 2001.
  CTAN: \url{macros/latex/contrib/ms/everyshi.dtx}
\end{thebibliography}

\loggingall
\end{document}
\endinput
%%
%% End of file `leaflet-manual.tex'.
